\section{Conclusion \& Limitation}
\label{sec:conclusion}
In this paper, we introduced \cubic, a training-free dynamic caching framework for efficient Flow Matching inference. By grounding the caching mechanism in the geometric properties of the velocity field—specifically, local trajectory curvature—\cubic adaptively partitions latents and optimizes update intervals to provably minimize error bounds. Comprehensive evaluations across diverse generation tasks demonstrate that \cubic establishes a superior quality-efficiency Pareto frontier compared to existing heuristic baselines. 
Nevertheless, our theoretical analysis is based on a forward-Euler abstraction of the sampling trajectory. While this offers a tractable link between curvature and cache-induced error, practical diffusion and flow-matching samplers often use multi-step, higher-order, or predictor--corrector updates, whose error constants and propagation patterns may differ. Therefore, our curvature-based scheduling rule should be viewed as a principled local approximation rather than a solver-agnostic guarantee, and extending the analysis to more general samplers is left for future work.

\section{Acknowledgements}
We would like to thank the anonymous reviewers for their insightful comments. This work is supported by the Fundamental and Interdisciplinary Disciplines Breakthrough Plan of the Ministry of Education of China (JYB2025XDXM910), NSFC for Distinguished Young Scholar under Grant 62225206, National Natural Science Foundation of China under Grants 62532006, 62495062, U23A600, and Beijing Natural Science Foundation under Grants L242017, L243001, L2603007. Jidong~Zhai is the corresponding author of this paper.



% We presented \cubic, a curvature-driven caching framework that translates ODE trajectory geometry into provably optimal cache scheduling for Flow Matching. Evaluations across diverse tasks confirm its superior quality-efficiency trade-offs over prior methods. Currently, the inherent memory overhead of layer-wise KV caches limits deployment on memory-constrained devices; future work will focus on frozen KV compression and efficient CPU offloading to mitigate this bottleneck.